% answers_short.tex —— 1977 年·江苏连云港市初试卷 · 简短答案速查
% 由 tools/make_short_answers.py 自 answers.tex 的 \gksolhead/\gkanswer 提取，
% 并与 1977/1977_数学真题_江苏连云港市初试_答案要点.md 逐题比对后生成，勿手改（改源 MD 后重跑）。

\begin{gknotice}
\textbf{使用说明}：本节只给最终答案便于快速核对，完整推导见解析版。答案为编写者独立推导，\textbf{非官方标准答案}。
\end{gknotice}

\gkpart{题目}
\begin{itemize}[leftmargin=2.6em]
  \item[\textbf{1.}] (1) $\log_3 9=2$；(2) ① $a_n=2n$；② $a_n=(-1)^n$；(3) $x^2+y^2=16$
  \item[\textbf{2.}] (1) $\dfrac{8}{27}$；(2) $-\dfrac23$；(3) $x^{\frac{13}{15}}$（$x>0$）；(4) $2+\sqrt2$
  \item[\textbf{3.}] (1) $x\in\mathbf R$ 且 $x\ne\dfrac12$；(2) $x=0$；(3) $x\ge1$
  \item[\textbf{4.}] $-2$。
  \item[\textbf{5.}] (1) $k=-\dfrac74$；(2) 当 $a\ne0$ 时，判别式 $\Delta=(a+b)^2+4a^2>0$ 恒成立，方程有两个不相等的实数根；当 $a=0$、$b\ne0$ 时，方程退化为一次方程，有唯一实根 $x=-1$；当 $a=b=0$ 时，方程对任意 $x$ 都成立
  \item[\textbf{6.}] (1) 长轴长 $10$，离心率 $e=\dfrac45$，焦点坐标 $(\pm4,0)$，准线方程 $x=\pm\dfrac{25}{4}$；(2) $A\left(4,\dfrac95\right)$；(3) $\dfrac{18}{5}$
  \item[\textbf{7.}] $BE=5$ 时扩充部分的面积最大，最大面积为 $41$。
  \item[\textbf{8.}] $BC=2\sqrt{13}$。
\end{itemize}
